% sections/k2k.tex — Kernel-to-Kernel Messaging
\section{Kernel-to-Kernel Messaging}
\label{sec:k2k}

The actor model treats messages as the only means of inter-actor
communication.  Translating that primitive to a GPU is one of the
principal implementation questions of the paper: a CPU actor system
can lean on operating-system mailboxes (heap-allocated,
lock-protected, accessed via syscalls); a GPU actor cannot.  A GPU actor lives in a kernel
that is itself a thread of a thread block of a cluster of a grid, with
distinct memory tiers (registers \(\to\) shared \(\to\) DSMEM \(\to\)
L2 \(\to\) HBM) and distinct synchronization primitives at each tier.
Choosing where to put the mailbox determines the latency and
bandwidth of every message the actor will ever send.

This section defines the kernel-to-kernel (K2K) messaging calculus.
We give a uniform syntactic envelope that abstracts the channel tier,
a small-step semantics for sending and delivering messages, and the
backpressure and dead-letter discipline that the runtime layers on
top.  We then prove three safety properties that the underlying
queue semantics implies: no-loss, FIFO per (source,\,destination),
and—once the idempotency cache is in place—at-most-once observable
delivery.  All three are mechanized in
\texttt{docs/verification/k2k\_delivery.tla}.

\paragraph{A note on delivery guarantees}
Three delivery properties show up below, and they are easy to
confuse.  We use them with the following specific meanings, once
and for all.  \textbf{No-loss} is an accounting property over the
envelope once it has been accepted by \textsc{SendOK}: the envelope
is either in flight, delivered, or dead-lettered with an auditable
reason—it never disappears (\Cref{thm:no-message-loss}).
\textbf{At-least-once} is an operational property: the runtime
exposes a replay API so the operator can re-inject backpressure
dead-letters until they deliver; at-least-once is therefore a
property of the calculus plus an explicit operator commitment
(\Cref{sec:k2k:dlq}).  \textbf{At-most-once observable delivery} is
the calculus property that the idempotency cache gives us: under
the cache-sizing hypothesis of
\Cref{thm:no-duplicate-delivery}, no two envelopes in \(D\) share
\((\mathit{mid}, \mathit{src})\); duplicates are short-circuited to
\(X\) with reason \state{IdempotencyReject} rather than admitted
to the visible delivery set.  The three properties are
complementary: no-loss gives accounting, at-least-once gives
operational recovery, at-most-once observable gives the visible
delivery cardinality used by audit and by downstream application
logic.  A typical trace over an accepted envelope is
\(\textsc{SendOK} \to \textsc{Deliver}\) (the happy path);
\(\textsc{SendOK} \to \textsc{DeliverDup} \to X\) if a duplicate
arrives after the first has been delivered (at-most-once
observable); and \(\textsc{SendFullDropOldest} \to X
\to \textsc{Replay} \to \textsc{SendOK} \to \textsc{Deliver}\) if
backpressure evicts the envelope and the operator chooses to
replay it (at-least-once via operator action).

\subsection{Channel Hierarchy}
\label{sec:k2k:tiers}

\system{} exposes three channel tiers, distinguished by which memory
substrate carries the message bytes:

\begin{enumerate}
\item \textbf{Intra-block (\textsc{Smem})}: source and destination
actors share a thread block; the queue is allocated in shared memory
(SMEM); enqueue/dequeue is a single \texttt{atomicAdd} on a 32-bit
producer/consumer cursor pair.  Latency: a handful of cycles.

\item \textbf{Intra-cluster (\textsc{Dsmem})}: source and destination
actors are in different blocks of the same cluster; the queue lives
in distributed shared memory (DSMEM); a remote-block pointer
materialized via \texttt{cluster.map\_shared\_rank()} feeds a
\texttt{cuda::barrier::arrive\_and\_wait} to publish the slot.
Latency: \(\sim\)30 cycles end-to-end on H100~\citep{nvidia2022hopper}.

\item \textbf{Inter-cluster (\textsc{Hbm})}: source and destination
actors are in different clusters; the queue lives in global memory
(HBM); enqueue is a \texttt{cuda::atomic\_ref<...>::fetch\_add} with
\texttt{memory\_order\_acq\_rel} ordering on the producer cursor.
Latency: an HBM round-trip ($\sim$400~cycles) plus L2 traversal.
\end{enumerate}

The calculus below abstracts the tier; we write \(\tau \in
\{\textsc{Smem}, \textsc{Dsmem}, \textsc{Hbm}\}\) for the
channel-tier tag carried in the envelope, and the small-step rules
treat all three uniformly.  The runtime makes the placement decision
(\Cref{sec:implementation}); the calculus is concerned only with the
end-to-end semantics each tier implements.
\Cref{fig:k2k-tiers} sketches the three tiers and their relationship
to the GPU memory hierarchy.

\begin{figure}[htbp]
\centering
\begin{tikzpicture}[
    every node/.style={font=\footnotesize\sffamily},
    block/.style={draw, rounded corners=1pt, minimum width=14mm,
                  minimum height=6mm, inner sep=2pt, font=\scriptsize\sffamily},
    cluster/.style={draw, dashed, rounded corners=3pt, inner sep=4mm,
                    fill=blue!4},
    gpu/.style={draw, very thick, rounded corners=4pt, inner sep=6mm,
                fill=gray!4},
    arrow/.style={->, >=Stealth, semithick}
  ]
  \begin{scope}
  \node[block] (b1) at (0, 0)   {Block 1};
  \node[block] (b2) at (1.7, 0) {Block 2};
  \node[block] (b3) at (3.4, 0) {Block 3};
  \node[block] (b4) at (0, -1.0)   {Block 4};
  \node[block] (b5) at (1.7, -1.0) {Block 5};
  \node[block] (b6) at (3.4, -1.0) {Block 6};
  \begin{pgfonlayer}{background}
    \node[cluster, fit=(b1)(b2)(b3)(b4)(b5)(b6),
          label={[anchor=north, font=\scriptsize]above:Cluster (1 GPC)}] (cl) {};
  \end{pgfonlayer}

  \node[block] (b7) at (6.5, -0.5) {Block 7};
  \node[block] (b8) at (8.2, -0.5) {Block 8};
  \begin{pgfonlayer}{background}
    \node[cluster, fit=(b7)(b8),
          label={[anchor=north, font=\scriptsize]above:Other Cluster}] (cl2) {};
  \end{pgfonlayer}
  \end{scope}

  \begin{pgfonlayer}{background}
    \node[gpu, fit=(cl)(cl2),
          label={[anchor=north west, font=\scriptsize]north west:GPU device (HBM)}] (g) {};
  \end{pgfonlayer}

  % Arrows
  \path[arrow, blue] (b1.east) edge[bend left=10]
        node[above, font=\tiny]{\textsc{Smem} (intra-block, $\sim$0 cy)} (b1.east);
  \draw[arrow, blue, dashed] (b1) to[bend left=18]
        node[above, font=\tiny]{\textsc{Dsmem} ($\sim$30 cy)} (b3);
  \draw[arrow, red, thick] (b6.east) to[bend right=20]
        node[below, font=\tiny]{\textsc{Hbm} ($\sim$400 cy)} (b7.west);

\end{tikzpicture}
\caption{K2K channel tiers.  \textsc{Smem} (intra-block) lives in
each block's shared memory; \textsc{Dsmem} (intra-cluster) uses
distributed shared memory addressed via
\texttt{cluster.map\_shared\_rank()}; \textsc{Hbm} (inter-cluster)
goes through global memory.  Annotated latencies are NVIDIA's
documented \emph{device-cycle} costs per
access~\citep{nvidia2022hopper}; the host-observed wall-clock for a
full kernel launch + synchronization round-trip on each tier is
reported separately in \Cref{tab:k2k-tier}.  The runtime selects
the tier based on the placement of source and destination actors.}
\label{fig:k2k-tiers}
\end{figure}

\subsection{Message Envelope}
\label{sec:k2k:envelope}

Every message transmitted on a K2K channel is wrapped in an
envelope—a fixed-layout header followed by a typed payload.  The
header is the unit on which routing, deduplication, dead-letter
disposition, and tracing operate; the payload is opaque to the K2K
layer.

\begin{definition}[K2K Envelope]
\label{def:k2k-envelope}
A K2K envelope is a tuple
\[
  e \;=\;
  \langle \mathit{src},\,
          \mathit{dst},\,
          \mathit{mid},\,
          \mathit{tag},\,
          \mathit{ts},\,
          \mathit{tier},\,
          \mathit{ttl},\,
          \mathit{payload} \rangle
\]
where \(\mathit{src}, \mathit{dst} \in \mathcal{A}\) are the source
and destination actor identifiers (\Cref{def:configuration});
\(\mathit{mid} \in \N\) is a globally unique message identifier
issued by the source;
\(\mathit{tag} \in \N \times \N\) is the audit boundary
\((\mathit{org\_id}, \mathit{engagement\_id})\) under which the
message is sent (\Cref{sec:tenancy});
\(\mathit{ts}\) is the source actor's HLC timestamp at the moment of
\textsc{Send} (\Cref{sec:supervision});
\(\mathit{tier} \in \{\textsc{Smem}, \textsc{Dsmem}, \textsc{Hbm}\}\)
is the chosen channel tier;
\(\mathit{ttl} \in \N\) is a hop counter decremented on every routing
step (\(\mathit{ttl} = 0 \Rightarrow\) drop to dead-letter); and
\(\mathit{payload}\) is the user-supplied byte sequence.  Let
\(\mathcal{E}\) denote the set of all envelopes.
\end{definition}

The choice of \(\mathit{mid} \in \N\) (rather than a hash of the
payload) is intentional: a single source may legitimately send
identical payloads in succession, and the calculus must treat them
as distinct.  The runtime allocates \(\mathit{mid}\) from a 64-bit
counter local to each source actor and tags the high bits with the
source identifier so identifiers are globally unique without
coordination.

\subsection{Configurations and Calculus}
\label{sec:k2k:calculus}

We extend the system configuration of \Cref{def:configuration} to
expose the K2K message store \(M\) explicitly.  In this section,
\(M\) is the disjoint union of three components:
\[
  M \;=\;
  Q \;\uplus\; D \;\uplus\; X
\]
where \(Q : \mathcal{A} \to \mathcal{E}^*\) maps each actor to its
inbox queue (a finite sequence of envelopes), \(D \subseteq
\mathcal{E}\) is the set of envelopes that have been delivered (i.e.,
consumed by a \textsc{Process} step), and \(X \subseteq \mathcal{E}
\times \Lambda\) is the dead-letter store, a set of envelope/reason
pairs where \(\Lambda\) is the set of dead-letter reasons defined in
\Cref{sec:k2k:dlq}.  We also assume an idempotency cache
\(I \subseteq \N \times \mathcal{A}\), recording the
\((\mathit{mid}, \mathit{src})\) pairs the runtime has already
accepted at least once at the destination side.

The K2K calculus consists of four small-step transitions:
\textsc{Send}, \textsc{Deliver}, \textsc{Drop}, and \textsc{Replay}.

\paragraph{Send}
The source attempts to enqueue an envelope \(e\) into the destination
queue \(Q(\mathit{dst})\).  \textsc{SendOK} fires only when the
destination is allocated and in \Active{} state (so its inbox is
accepting traffic), the source and destination are distinct, and
the queue has spare capacity.  If the destination is in
\Draining{}, \Initializing{}, \Terminated{}, or \Failed{}, a
separate \textsc{DropInvalidTarget} rule short-circuits to
dead-letter with reason \state{InvalidTarget}
(\Cref{sec:k2k:dlq}); if the queue is full, the configured
backpressure policy (\Cref{sec:k2k:backpressure}) applies.  We
give the success rule here; the two backpressure-failure rules
follow in \Cref{sec:k2k:backpressure}.

\begin{equation}
\inference[\textsc{SendOK}]
{
  e = \langle s, d, m, \tau, \ldots\rangle \\
  \Theta(d) = \langle \Active{}, \ldots\rangle \\
  |Q(d)| < \mathit{cap} \\
  d \neq s
}
{
  \langle \Theta, \langle Q, D, X\rangle, t\rangle
  \trans{\mathsf{send}(e)}
  \langle \Theta, \langle Q', D, X\rangle, t+1\rangle
}
\label{rule:k2k-send}
\end{equation}

\noindent where \(Q' = Q[d \mapsto Q(d) \cdot e]\) (i.e., \(e\)
appended).  The \(\Theta(d) = \langle \Active{}, \ldots\rangle\)
premise is the \emph{destination-lifecycle precondition} that the
appendix Quiesce rule (\Cref{rule:quiesce}) relies on to justify
rejecting sends to a \Draining{} actor.

\paragraph{Deliver}
A destination actor \(d\) consumes the head of its queue.  This
corresponds to the \textsc{Process} rule of \Cref{rule:process}
threading through one envelope; we expose it here as a separate K2K
event for traceability.

\begin{equation}
\inference[\textsc{Deliver}]
{
  Q(d) = e \cdot q \\
  e = \langle s, d, m, \tau, \ldots\rangle \\
  (m, s) \notin I
}
{
  \langle \Theta, \langle Q, D, X\rangle, t\rangle
  \trans{\mathsf{deliver}(e)}
  \langle \Theta, \langle Q', D', X\rangle, t+1\rangle
}
\label{rule:k2k-deliver}
\end{equation}

\noindent where \(Q' = Q[d \mapsto q]\), \(D' = D \cup \{e\}\), and
the runtime adds \((m, s)\) to \(I\) as a side effect.  The
\((m, s) \notin I\) precondition implements the at-most-once side of
the idempotency guarantee (\Cref{sec:k2k:idempotency}); a duplicate
envelope (\(\mathit{mid}, \mathit{src}\) already seen) is consumed
without emitting a delivery, via \textsc{DeliverDup} below.

\begin{equation}
\inference[\textsc{DeliverDup}]
{
  Q(d) = e \cdot q \\
  e = \langle s, d, m, \tau, \ldots\rangle \\
  (m, s) \in I
}
{
  \langle \Theta, \langle Q, D, X\rangle, t\rangle
  \trans{\mathsf{drop\_dup}(e)}
  \langle \Theta, \langle Q', D, X'\rangle, t+1\rangle
}
\label{rule:k2k-deliver-dup}
\end{equation}

\noindent where \(Q' = Q[d \mapsto q]\) and \(X' = X \cup
\{(e, \state{IdempotencyReject})\}\).  The duplicate is recorded in
the dead-letter store with the reason \state{IdempotencyReject}; this
ensures the delivery cardinality on the visible side
(\(D\)) stays bounded, while preserving full audit trace on the
dead-letter side (\(X\)).

\paragraph{Drop and Replay}
\textsc{Drop} (with reason) and \textsc{Replay} are detailed in
\Cref{sec:k2k:backpressure,sec:k2k:dlq}.

\subsection{Backpressure Policies}
\label{sec:k2k:backpressure}

When the destination queue is full at the moment of \textsc{Send}, the
runtime must choose a disposition.  \system{} supports two policies,
selectable per channel tier and per audit boundary:

\begin{description}
\item[\(\mathsf{RejectAndDLQ}\)] The envelope is not enqueued; it is
moved to the dead-letter store with reason
\state{Overflow}.  The source's \textsc{Send} call returns a failed
status which the source may observe and retry.

\item[\(\mathsf{DropOldest}\)] The envelope at the head of the
destination queue is moved to the dead-letter store with reason
\state{Evicted}, and the new envelope is enqueued at the tail.  This
prefers fresh data over stale; it tends to be the preferred policy
for telemetry and best-effort streams.
\end{description}

\begin{equation}
\inference[\textsc{SendFullReject}]
{
  e = \langle s, d, m, \tau, \ldots\rangle \\
  |Q(d)| = \mathit{cap} \\
  \mathit{policy}(d) = \mathsf{RejectAndDLQ}
}
{
  \langle \Theta, \langle Q, D, X\rangle, t\rangle
  \trans{\mathsf{drop}(e, \state{Overflow})}
  \langle \Theta, \langle Q, D, X'\rangle, t+1\rangle
}
\label{rule:k2k-send-reject}
\end{equation}

\noindent where \(X' = X \cup \{(e, \state{Overflow})\}\).

\begin{equation}
\inference[\textsc{SendFullDropOldest}]
{
  e = \langle s, d, m, \tau, \ldots\rangle \\
  Q(d) = e_0 \cdot q,\; |Q(d)| = \mathit{cap} \\
  \mathit{policy}(d) = \mathsf{DropOldest}
}
{
  \langle \Theta, \langle Q, D, X\rangle, t\rangle
  \trans{\mathsf{evict\_drop}(e_0)}
  \langle \Theta, \langle Q', D, X'\rangle, t+1\rangle
}
\label{rule:k2k-send-evict}
\end{equation}

\noindent where \(Q' = Q[d \mapsto q \cdot e]\) and
\(X' = X \cup \{(e_0, \state{Evicted})\}\).

\subsection{Idempotency Cache}
\label{sec:k2k:idempotency}

The cache \(I \subseteq \N \times \mathcal{A}\) records the
\((\mathit{mid}, \mathit{src})\) pairs that have been accepted at the
destination side, regardless of whether the source eventually retries.
Two operations interact with the cache:

\begin{itemize}
\item \textsc{Deliver}~(\Cref{rule:k2k-deliver}) requires
\((\mathit{mid}, \mathit{src}) \notin I\) as a precondition and adds
the pair to \(I\) as a side effect.
\item \textsc{DeliverDup}~(\Cref{rule:k2k-deliver-dup}) is the
contrapositive: the pair is already in \(I\), so the envelope is
moved to dead-letter rather than the visible delivery set.
\end{itemize}

The cache is bounded; the runtime evicts entries by an LRU policy with
configurable maximum size.  An evicted entry can in principle re-enter
\(I\) on a later sighting, which would re-admit the duplicate to the
visible delivery set and violate at-most-once.  In practice the cache
is sized so that the eviction window exceeds the maximum sender retry
window, an assumption we make explicit
(\Cref{thm:no-duplicate-delivery}, condition \emph{(c)}).

\subsection{Dead-Letter Discipline}
\label{sec:k2k:dlq}

The dead-letter store \(X\) records every envelope that left the
queue without ever crossing into \(D\).  The reason set is:
\[
  \Lambda \;=\;
  \{\state{Overflow},\,
    \state{Evicted},\,
    \state{InvalidTarget},\,
    \state{IdempotencyReject},\,
    \state{TenantMismatch},\,
    \state{TTLExhausted}\}.
\]

The first two are produced by the backpressure rules of
\Cref{sec:k2k:backpressure}; \state{IdempotencyReject} by
\Cref{rule:k2k-deliver-dup}; \state{InvalidTarget} when
\(d \notin \dom(\Theta)\) at the time of routing;
\state{TenantMismatch} when the source's tag and the destination's
tag would cross a tenant boundary (\Cref{sec:tenancy}); and
\state{TTLExhausted} when the routing layer decremented
\(\mathit{ttl}\) to zero before reaching \(\mathit{dst}\).

The discipline is total: every envelope that does not reach \(D\)
is in \(X\) with a reason.  The system does not drop a message
without recording it.  This property underwrites the audit
guarantee of \Cref{sec:tenancy:audit} and also lets the delivery
theorems in \Cref{sec:k2k:theorems} be stated as a conjunction of
trace properties over \((D, X)\).

\paragraph{Replay}
The dead-letter store can be inspected and selectively re-injected by
a supervisor (the runtime exposes
\texttt{POST /api/v1/admin/k2k/dlq/replay}).  We model this as a
single rule:

\begin{equation}
\inference[\textsc{Replay}]
{
  (e, \lambda) \in X \\
  \lambda \in \{\state{Overflow}, \state{Evicted}\} \\
  e = \langle s, d, m, \tau, \ldots\rangle \\
  d \in \dom(\Theta)
}
{
  \langle \Theta, \langle Q, D, X\rangle, t\rangle
  \trans{\mathsf{replay}(e)}
  \langle \Theta, \langle Q, D, X\rangle, t+1\rangle
}
\label{rule:k2k-replay}
\end{equation}

\noindent followed by an immediate \textsc{Send}(\(e\)) using the
fresh state (or a fresh \textsc{SendFullReject} if the queue is
again full).  Replay only applies to backpressure reasons
(\state{Overflow}, \state{Evicted}); replaying an envelope
dropped for \state{InvalidTarget}, \state{TenantMismatch},
\state{IdempotencyReject}, or \state{TTLExhausted} would re-trigger
the same drop and is forbidden by the rule's premise.

\subsection{Safety Theorems}
\label{sec:k2k:theorems}

We can now state the three safety properties that the calculus
guarantees.  Each corresponds directly to an invariant in the
mechanized TLA\ensuremath{^{+}} spec.

\begin{theorem}[No Message Loss]
\label{thm:no-message-loss}
For every reachable configuration \(\langle \Theta,
\langle Q, D, X\rangle, t\rangle\) and every envelope \(e\) for which
a \textsc{SendOK}~(\Cref{rule:k2k-send}) step has been applied along
the trace, exactly one of the following holds:
\begin{enumerate}
\item \(e \in D\) (delivered to the visible delivery set by
\textsc{Deliver});
\item \(e \in Q(\mathit{dst})\) at some position (in flight);
\item \((e, \lambda) \in X\) for some
\(\lambda \in \{\state{Evicted},\, \state{IdempotencyReject}\}\)
(moved to dead-letter either because backpressure evicted it under
the \(\mathsf{DropOldest}\) policy or because the idempotency cache
rejected it as a duplicate).
\end{enumerate}
The three alternatives are mutually exclusive: an envelope cannot
simultaneously be in \(D\), in some \(Q(\mathit{dst})\), and in
\(X\).  Together they cover every envelope that \textsc{SendOK}
has accepted; nothing is silently dropped.  Dead-letter disposition
for the other reasons in \(\Lambda\) (\state{Overflow},
\state{InvalidTarget}, \state{TenantMismatch}, \state{TTLExhausted})
applies to envelopes that \textsc{SendOK} never accepted and is
therefore outside the scope of this theorem.
\end{theorem}

\begin{proof}[Proof sketch]
Induction on the trace.  The base case is vacuous (no \textsc{Send}
has been applied).  For the induction step, consider the last rule
applied:
\begin{itemize}
\item \textsc{SendOK}~(\Cref{rule:k2k-send}) appends \(e\) to
\(Q(\mathit{dst})\), so condition~(2) holds for \(e\) immediately
after the step; conditions~(1) and~(3) are unaffected.
\item \textsc{Deliver}~(\Cref{rule:k2k-deliver}) removes the head of
\(Q(\mathit{dst})\) and adds it to \(D\), so the disposition of that
envelope moves from~(2) to~(1); other envelopes in flight retain
their disposition.
\item \textsc{DeliverDup}~(\Cref{rule:k2k-deliver-dup}) removes the
head of \(Q(\mathit{dst})\) and adds \((e, \state{IdempotencyReject})\)
to \(X\), so the envelope moves from~(2) to~(3).
\item \textsc{SendFullDropOldest}~(\Cref{rule:k2k-send-evict})
removes the head of \(Q(\mathit{dst})\) and adds \((e_0,
\state{Evicted})\) to \(X\), so the evicted envelope's disposition
moves from~(2) to~(3).
\end{itemize}
Mutual exclusion holds because no rule places an envelope into one
of (1)--(3) without first removing it from another, and the three
targets (\(D\), \(Q\), \(X\)) are set-disjoint by construction.  By
induction, the partition ``\(e \in D\) \emph{xor} \(e \in
Q(\mathit{dst})\) \emph{xor} \((e, \state{Evicted}) \in X\) \emph{xor}
\((e, \state{IdempotencyReject}) \in X\)'' is preserved at every
step.
\qed
\end{proof}

\Cref{thm:no-message-loss} is the analogue of the
\texttt{NoMessageLoss} invariant in
\texttt{docs/\allowbreak verification/\allowbreak
k2k\_delivery.tla}, lines 95--104.  The
mechanization elides the dead-letter store and treats only the
in-flight/delivered case (clauses~(1) and~(2)), which is an
appropriate abstraction for a small bounded model; the analytical
proof above covers the full discipline, including dead-letter
dispositions.

\begin{theorem}[FIFO per Source--Destination Pair]
\label{thm:fifo-per-sender}
For all reachable configurations and for every pair of envelopes
\(e_1 = \langle s, d, m_1, \ldots\rangle\) and
\(e_2 = \langle s, d, m_2, \ldots\rangle\) sharing the same source
\(s\) and destination \(d\) such that \textsc{SendOK}\((e_1)\) was
applied before \textsc{SendOK}\((e_2)\) (i.e., \(m_1 < m_2\) since
the source assigns identifiers monotonically), if both \(e_1\) and
\(e_2\) appear in \(D\), then \textsc{Deliver}\((e_1)\) was applied
before \textsc{Deliver}\((e_2)\).
\end{theorem}

\begin{proof}[Proof sketch]
The destination queue \(Q(d)\) is a sequence; \textsc{SendOK}
appends to the tail and \textsc{Deliver} consumes the head.  Both
operations preserve the relative order of the remaining elements.
By induction on the trace, the relative order of any two envelopes
that share \((s, d)\) and are both alive in \(Q(d)\) at any point
is the order in which they were appended, which is the order in
which the source assigned their \(\mathit{mid}\)s.  Therefore, if
\(\mathit{mid}_1 < \mathit{mid}_2\), then \(e_1\) is delivered no
later than \(e_2\).  The presence of \textsc{SendFullDropOldest}
between the two does not affect the property: the only envelope
\textsc{SendFullDropOldest} can evict is the head of the queue, and
if it evicts \(e_1\) before \(e_1\) is delivered, then \(e_1
\notin D\) and the antecedent of the implication fails.
\qed
\end{proof}

\Cref{thm:fifo-per-sender} corresponds to \texttt{FIFOPerSender} in
the mechanized spec (lines 108--115).  Note that the property is
stated only for envelopes that share \((\mathit{src}, \mathit{dst})\):
two envelopes from the same source to different destinations may
arrive in either order, and two envelopes from different sources to
the same destination interleave according to the queue's append
schedule.

\begin{theorem}[At-Most-Once Observable Delivery]
\label{thm:no-duplicate-delivery}
Suppose the idempotency cache is sized so that no
\((\mathit{mid}, \mathit{src})\) pair is evicted from \(I\) within
the source's retry window for that envelope.  Then for every
reachable configuration, no envelope identifier \((\mathit{mid},
\mathit{src})\) appears in two distinct envelopes in \(D\).
\end{theorem}

\begin{proof}[Proof sketch]
Suppose for contradiction that \(D\) contains two envelopes
\(e_1 \neq e_2\) with \(e_1.\mathit{mid} = e_2.\mathit{mid}\) and
\(e_1.\mathit{src} = e_2.\mathit{src}\).  Both must have been
admitted to \(D\) by \textsc{Deliver}~(\Cref{rule:k2k-deliver}),
which has the precondition
\((\mathit{mid}, \mathit{src}) \notin I\) and the side effect of
adding \((\mathit{mid}, \mathit{src})\) to \(I\).  Without loss of
generality say \textsc{Deliver}\((e_1)\) was applied first.  After
that step, \((\mathit{mid}, \mathit{src}) \in I\).  By the
hypothesis on cache sizing, \((\mathit{mid}, \mathit{src})\)
remains in \(I\) until the source's retry window for that envelope
has closed; the source assigns a fresh \(\mathit{mid}\) for any
subsequent envelope.  But \(e_2\) shares \(\mathit{mid}\) with
\(e_1\), so \(e_2\) was sent within the retry window, contradicting
the precondition of \textsc{Deliver}\((e_2)\).
\qed
\end{proof}

\Cref{thm:no-duplicate-delivery} corresponds to
\texttt{NoDuplicateDelivery} in the mechanized spec (lines 121--123),
which states the same property in the simpler form ``no two distinct
delivery tuples share an \(\mathit{mid}\)'' (the spec elides the
\(\mathit{src}\) component because the spec's \(\mathit{mid}\)s are
globally unique by construction).  The
analytical statement above adds the cache-sizing precondition
explicitly because in the runtime the cache is finite, and the
hypothesis is something the operator must verify against the actual
retry policy.

\paragraph{Implementation note}
A weaker theorem—at-least-once delivery—follows immediately from
\Cref{thm:no-message-loss} together with the operator's commitment
to invoke \textsc{Replay} on every \state{Overflow} dead-letter.
We do not state this as a separate theorem because at-least-once is
not a property the calculus alone can guarantee: it depends on
operator behavior outside the formal model.  The runtime exposes
the dead-letter API for exactly this purpose
(\Cref{sec:implementation:admin}).

\subsection{Mechanization}
\label{sec:k2k:mechanization}

The K2K calculus is mechanized in \TLA{} in
\texttt{docs/verification/k2k\_delivery.tla}.  The mechanization
keeps four state variables (\texttt{queues}, \texttt{sent\_accepted},
\texttt{delivered}, \texttt{dropped}) and two actions (\texttt{Send},
\texttt{Deliver}); it elides the dead-letter reason set, the
idempotency cache, and the replay rule, modeling them as sources of
nondeterminism.  We model-checked under the configuration
\texttt{k2k\_delivery.cfg} for parameter ranges
\(|\textit{Actors}| \le 3\), \(\textit{MaxMsgs} \le 4\), and
\(\textit{QueueCap} \le 2\); within those bounds the invariants
\texttt{NoMessageLoss}, \texttt{FIFOPerSender},
\texttt{NoDuplicateDelivery}, \texttt{DropAcceptedDisjoint}, and
\texttt{QueueBoundOK} hold on every reachable state, with the search
exhausting the explored state space without violation.
